\documentclass{elsarticle}
\usepackage{lineno}  

\usepackage{titletoc}
\usepackage[T1]{fontenc}
\usepackage{etoc}
\usepackage[utf8]{inputenc}
\usepackage{lmodern}
\usepackage{pythonhighlight}
\usepackage{xcolor}
\usepackage{url}
\usepackage{rotating}
\usepackage{setspace}
\usepackage{mathrsfs}
%\usepackage{natbib}
\usepackage[english]{babel}
\usepackage{multirow}
\usepackage[hidelinks]{hyperref}
\usepackage{graphicx}
\usepackage{caption}
\usepackage{subcaption}
\usepackage{amsmath,amsfonts,amssymb,amsthm}
\usepackage{enumerate}
\usepackage{geometry}
\usepackage{mathtools}
\usepackage{enumitem}
\usepackage{cases}
\usepackage{mathabx}
\usepackage{environ}
\usepackage{epstopdf}
\usepackage{cleveref}
\usepackage{xspace}
\usepackage{float}
\usepackage{comment}
\usepackage{pdflscape}
\usepackage{enumitem}
\usepackage{graphicx}
\usepackage{color}
\usepackage{subcaption}

\usepackage{amsmath,amsfonts,amssymb,dsfont}
\usepackage{siunitx}
\usepackage{booktabs}
\usepackage{amsopn}
\usepackage[english]{babel}
\usepackage[autostyle]{csquotes}
\usepackage{makeidx}

\DeclareMathOperator{\diag}{diag}
\DeclareMathOperator{\rank}{rank}
\newcommand{\Ai}{\mathscr{A}_1}
\newcommand{\Aii}{\mathscr{A}_2}
\newcommand{\Ae}{\mathscr{A}_i}
\newcommand{\As}{\mathscr{A}}
\newcommand{\Rc}{\mathcal{R}}
\newcommand{\R}{\mathbb{R}}
\newcommand{\Nc}{\mathcal{N}}
\newcommand{\Ec}{\mathcal{E}}
\newcommand{\ds}{\displaystyle}
\def\G{\mathcal{G}}
\def\L{\mathcal{L}}
\def\M{\mathcal{M}}
\def\P{\mathcal{P}}
\def\R{\mathcal{R}}
\def\U{\mathcal{U}}
\def\M{\mathcal{M}}
\def\I{\mathcal{I}}
\def\F{\mathcal{F}}
\def\V{\mathcal{V}}
\def\bb{\mathbf{b}}
\def\bf{\mathbf{f}}
\def\bh{\mathbf{h}}
\def\bx{\mathbf{x}}
\def\bB{\mathbf{B}}
\def\bD{\mathbf{D}}
\def\bE{\mathbf{E}}
\def\bI{\mathbf{I}}
\def\bN{\mathbf{N}}
\def\bR{\mathbf{R}}
\def\bS{\mathbf{S}}
\def\bV{\mathbf{V}}
\def\bW{\mathbf{W}}
\def\bGamma{\mathbf{\Gamma}}
\def\bx{\mathbf{x}}
\def\bPi{\mathbf{\Pi}}
\def\b0{\mathbf{0}}
\def\nbOne{{\mathchoice {\rm 1\mskip-4mu l} {\rm 1\mskip-4mu l}{\rm 1\mskip-4.5mu l} {\rm 1\mskip-5mu l}}}
\def\diag{\mathsf{diag}}

\usepackage{tikz}
\usetikzlibrary{shapes,arrows}
\usetikzlibrary{positioning}
\tikzstyle{cloud} = [draw, ellipse,fill=red!20, node distance=0.87cm,
minimum height=2em]
\tikzstyle{line} = [draw, -latex']
\usetikzlibrary{shapes.symbols,shapes.callouts,patterns}
\usetikzlibrary{calc}
\tikzstyle{myBlock}=[draw, minimum size=0.5in, node distance=1.5in]
\usetikzlibrary{arrows.meta}

\input xy
\xyoption{all}	

\usepackage{amsmath,amsfonts,amssymb,dsfont}

%%%%%%%%%%% Defining Enunciations  %%%%%%%%%%%
\newtheorem{theorem}{\bf Theorem}[section]
\newtheorem{condition}{\bf Condition}[section]
\newtheorem{corollary}{\bf Corollary}[section]
%%%%%%%%%%%%%%%%%%%%%%%%%%%%%%%%%%%%%%%%%%%%%%%

\begin{document}

\title{Temporal Network Restructuring of Air Travel During COVID-19: Implications for Dynamic Infectious Disease Surveillance}


\begin{abstract}
\end{abstract}

\maketitle

\section{Introduction}

Transport and health research has traditionally focused on active travel modes (walking and cycling) and their direct effects on physical activity, road safety, and local air quality. These are important determinants of health at the community scale. Global air transportation operates at an entirely different scale: it is the macro-level transport determinant of health that governs how local outbreaks become regional epidemics and regional epidemics become pandemics. The structural properties of air networks shape epidemic dynamics in ways that have no parallel in active travel research, and understanding those properties is a prerequisite for effective border health surveillance.

Surveillance systems for infectious disease importation rely primarily on monitoring passenger volumes at borders. During the COVID-19 pandemic, this approach revealed a fundamental limitation: passenger counts measure how many people are moving, but not the structure of the network through which they move. These are different things, and they carry different implications for importation risk. An infected traveler arriving at a major hub airport with connections to dozens of onward destinations poses a different risk than one arriving at an isolated regional airport, even if the passenger count is the same.

Prior work has established that air transportation network structure shapes epidemic dynamics in fundamental ways. Colizza et al.\ demonstrated that the heterogeneous structure of the global airline network, rather than simple passenger volume, determines the predictability and speed of global epidemic spread \cite{colizza2006}. Brockmann and Helbing showed that geographical distance is a poor predictor of disease arrival time, and that effective distance through transportation networks provides a more accurate measure of epidemic risk \cite{brockmann2013}. Despite this foundation, most operational disease surveillance systems do not incorporate network structure, and most network analyses of air transportation use static snapshots from a single year or aggregate data into yearly totals. This obscures precisely what public health practitioners need to know during a crisis: how the network is changing right now, and what those changes mean for importation risk this month versus last month.

The COVID-19 pandemic exposed this gap acutely. In March and April 2020, as travel bans took effect and airlines cancelled flights, the global air transportation network did not simply shrink: it restructured. Routes disappeared selectively. Communities of connected airports fragmented. Hub airports remained critical connectors even as the overall system contracted. A surveillance system calibrated only to passenger volume would have registered a 92\% reduction in Canadian arrivals in April 2020 and concluded that importation risk had dropped proportionally \cite{StatisticsCanada_2026}. A system accounting for network structure would have identified something more nuanced: that the remaining travel was increasingly concentrated through fewer, more critical pathways, and that the reconstruction of those pathways in 2021 was reopening importation routes faster than passenger volumes suggested.

We analysed air transportation data spanning January 2019 through December 2021, covering the pre-pandemic baseline, the collapse of 2020, and the partial recovery of 2021. Using freely available flight data from the OpenSky Network, we constructed 36 monthly directed networks covering airports in Canada, the United States, and Europe, and tracked structural network metrics across the full pandemic cycle. Rather than modelling disease spread directly, we characterize how the network itself changed and assess whether those structural changes are consistent with documented patterns of SARS-CoV-2 variant importation into Canada. 

The paper is organized as follows. Section~\ref{sec:methods} describes the data, processing steps, and network construction. Section~3 presents the results, beginning with the temporal alignment between network conditions and documented variant importation events, followed by the structural network metrics that explain the observed patterns. Section~4 discusses implications for border health surveillance and pandemic preparedness, and proposes concrete directions for integrating network intelligence into operational public health practice.

\section{Methods}\label{sec:methods}

\subsection{Data Source}\label{subsec:data_source}

We used flight data from the OpenSky Network, which collects information from aircraft broadcasts \cite{OpenSky_dataset}. Every airplane transmits its position, speed, and identification via a system called ADS-B (Automatic Dependent Surveillance-Broadcast). Anyone with a simple ground receiver can pick up these signals and share them with the OpenSky Network. This creates a crowdsourced database of flight information that is freely available for research.

The OpenSky data has several notable limitations. ADS-B broadcasts are mandated by law in some places (Canada, the USA, Europe) \cite{CAN_ADSB, US_ADSB, EU_ADSB} but not others. This means coverage is good in these regions but spotty elsewhere. If an aircraft does not fly in countries with ADS-B requirements, it probably will not be in the database. There are also some technical issues related to trajectory processing \cite{Baek_ADS-B}. For example, if a small airport sits close to a major hub, flights might get attributed to the wrong airport. Sometimes origin and destination data is missing entirely when the system cannot identify an airport. Passenger volumes can only be estimated when the aircraft type is matched to its capacity.

We used data processed by Xavier Olive et al., available on Zenodo \cite{xavier_olive_2022}. They aggregated the OpenSky raw data into monthly files covering January 2019 through March 2022. We focused on 2019, 2020, and 2021 (36 months total), capturing a full year of normal travel, the pandemic collapse, and the beginning of recovery. One important caveat applies to the 2021 figures: the OpenSky volunteer receiver network grew substantially between 2019 and 2021, with additional ground stations deployed across North America and Europe. Some portion of the observed increase in edge counts in 2021 above pre-pandemic levels may therefore reflect improved ADS-B coverage detecting flights that were also present in 2019 but not captured, rather than genuine network expansion. We cannot fully disentangle coverage growth from real route growth in the absence of a ground-truth flight database for comparison; the 2021 structural findings should accordingly be interpreted as an upper bound on true network expansion.

\subsection{Cleaning and Preparing the Data}

The data came pre-cleaned to some extent, but additional processing was required. We removed flights with missing origin or destination information and flights where the origin and destination were the same airport. To illustrate the scale of this processing: in January 2019, we started with approximately 2.66 million flight records and retained 1.32 million after cleaning. In April 2020, at the pandemic's lowest point, we had 842,905 initial records and retained 492,877. The collapse in April 2020 is visible in both raw and cleaned numbers.

We enriched the data by adding country and continent information for each airport using a standard airport database \cite{airport_code_dataset}, and aircraft capacity data \cite{aircraft_capacity} to estimate passenger volume per flight. For flights where the aircraft type could not be identified, we assigned a volume of 2 passengers, capturing the small private planes common in the United States that would otherwise be unaccounted for.

We filtered to airports in Canada, the United States, and Europe. Europe includes 44 countries: Albania, Andorra, Armenia, Austria, Belarus, Belgium, Bosnia and Herzegovina, Bulgaria, Croatia, Cyprus, Czech Republic, Denmark, Estonia, Finland, France, Germany, Greece, Hungary, Iceland, Ireland, Italy, Latvia, Liechtenstein, Lithuania, Luxembourg, Malta, Moldova, Monaco, Montenegro, Netherlands, North Macedonia, Norway, Poland, Portugal, Romania, Russia, San Marino, Serbia, Slovakia, Slovenia, Spain, Sweden, Switzerland, Turkey, Ukraine, and the United Kingdom. For US territories outside the continental US, we included only Puerto Rico; for European territories, only those proximate to the continent.

\subsection{Building the Network}

For each month, we built a directed network where airports are nodes and flights are edges. The weight on each edge represents the total estimated number of passengers travelling that route in that month. For example, if we observe 10 flights from Toronto to New York in January using aircraft with an average capacity of 200 seats, the edge from Toronto to New York receives a weight of approximately 2,000 passengers. This approach yields 36 consecutive monthly snapshots from January 2019 through December 2021.

\subsection{Computing Network Metrics}

We calculated several types of metrics for each monthly network. Network-level metrics characterize overall structure: node count (airports), edge count (routes), density, average clustering coefficient, diameter, and number of strongly connected components. Centrality measures identify structurally important airports: degree centrality, betweenness centrality, PageRank, harmonic centrality, eigenvector centrality, and eccentricity.

From a public health perspective, each centrality measure captures a distinct dimension of disease importation risk. \textit{Betweenness centrality} quantifies how often an airport lies on the shortest path between other airports; high-betweenness airports act as bridge points where a single infected traveler can seed multiple new regions simultaneously. \textit{PageRank} measures the quality of an airport's connections rather than their quantity; high-PageRank airports act as amplifiers because they are well-connected to other major hubs, accelerating the onward dispersal of any pathogen introduced there. \textit{Closeness centrality} measures the average number of hops required to reach all other nodes; a high closeness score means an infected traveler could reach any destination in the network quickly. \textit{Degree centrality} (in-degree and out-degree) captures raw connectivity: airports with many incoming routes face greater importation exposure, while airports with many outgoing routes have greater potential to seed downstream locations.

Community detection groups airports into clusters based on traffic patterns. We used two algorithms: Louvain and Infomap. Louvain optimizes for modularity, while Infomap identifies communities by minimizing the description length of a random walker's path. Running both algorithms allows us to assess whether community structure is stable or sensitive to the choice of algorithm. We tracked the number of communities and the size of the largest community over time.

Edge dynamics capture month-to-month network change: we counted edges appearing for the first time (new routes) and edges disappearing (cancelled routes).

\subsection{Why These Metrics Matter for Epidemiology}

Network metrics translate directly into disease risk. A high-betweenness airport acts as a bridge for disease spread; a high-PageRank airport is well-connected to other hubs, making it a potential amplifier. Community structure determines whether disease would spread globally or remain contained within regional clusters. Edge dissolution in April 2020 represents travel restrictions and route cancellations that directly reduced transmission pathways.

By tracking all metrics across 36 months, we can observe how the network adapted to the pandemic and measure precisely how the crisis changed the structure through which disease spread. This work directly supports epidemiological research aimed at estimating disease importation risk during the early stages of outbreaks, as outlined in our June 2020 report to the Public Health Agency of Canada \cite{PHAC2020}. The temporal network analysis framework developed here enables dynamic assessment of how travel restrictions impact disease propagation and identifies high-risk points of entry.


\section{Results}

\subsection{Temporal Alignment of Network Recovery with Documented Variant Importation into Canada}\label{subsec:variant_alignment}


The central question for border health surveillance is whether structural network conditions are consistent with observed patterns of disease importation. To address this, we examine the timing of three major SARS-CoV-2 variant importation events into Canada against the network evolution documented across 36 monthly snapshots from 2019 through 2021. We do not claim a causal relationship; rather, we assess whether documented variant arrivals are temporally consistent with the network conditions that would be expected to facilitate or impede importation. Direct causal validation would require travel-linked case surveillance data at a resolution not publicly available for Canada. The structural metrics supporting this analysis are presented in Sections~3.2 through~3.6.

\subsubsection*{Alpha (B.1.1.7)}

The Alpha variant originated in the United Kingdom in early autumn 2020 and was designated a variant of concern by the WHO in December 2020 \cite{rambaut2020preliminary}. Its first confirmed detection in Canada occurred in Ontario in late December 2020 \cite{alpha_canada_2020}.

This timing is consistent with the edge formation dynamics we observe. Figure~\ref{fig:ADSB_edges_formed} shows that route formation in 2020 reached its lowest point in April (approximately 79,000 new edges), before recovering through the summer and autumn. By October through December 2020, edge formation had recovered to approximately 130,000--150,000 routes per month, levels comparable to the equivalent months in 2019. Alpha's first detection in Canada in late December 2020 therefore falls within this recovery window: the routes through which it could travel had been largely restored by the time it established itself, but were absent or severely reduced during the April through September period when Alpha was emerging and circulating in the United Kingdom.

\subsubsection*{Delta (B.1.617.2)}

The Delta variant was first identified in India in late 2020 and designated a WHO variant of concern in May 2021 \cite{WHO_delta_2021}. It became the dominant strain circulating in Canada by summer 2021 \cite{PHAC_variants}.

Delta's importation into Canada falls squarely within the period our analysis identifies as active network reconstruction. Figure~\ref{fig:ADSB_edges_formed} shows that edge formation in 2021 was substantially elevated above both 2019 and 2020 throughout the year, with values consistently above 145,000--195,000 routes per month. This represents a network that was not merely recovering but actively adding routes at an accelerated rate. Notably, Canadian border arrivals remained well below pre-pandemic levels throughout this period, yet Delta still established itself nationally. This is consistent with the structural observation made in Section~\ref{subsec_recovery_Persistent}: network topology had recovered more rapidly than passenger volume. The routes through which the variant could travel had reopened even while total traveller numbers remained depressed, meaning that a structurally recovering network operating at partial passenger capacity still provided sufficient pathways for variant introduction and domestic spread.

\subsubsection*{Omicron (B.1.1.529)}

The Omicron case provides the most striking temporal alignment between network conditions and documented importation. Omicron was first reported to the WHO on November 24, 2021, and classified as a variant of concern on November 26, 2021 \cite{WHO_omicron_2021}. Canada's first confirmed cases were announced on November 28, 2021, in Ottawa, Ontario, in two individuals with recent travel history from Nigeria \cite{ontario_omicron_2021}. Additional cases were confirmed in British Columbia on November 30, and in Alberta and Quebec within days \cite{bc_omicron_2021}. Omicron reached Canada within four days of its global WHO classification.

This near-immediate arrival occurred precisely during the period our analysis identifies as the most advanced stage of network reintegration. By November 2021, Figure~\ref{fig:ADSB_edges_formed} shows edge formation running consistently at approximately 183,000 routes per month, substantially above both 2019 and 2020 levels for the same month. Edge dissolution in 2021 (Figure~\ref{fig:ADSB_edges_dissolved}) was similarly elevated throughout the year, reflecting high network turnover characteristic of active restructuring rather than the static compression of 2020. Community structure had reverted toward pre-pandemic configurations across both detection algorithms (Figure~\ref{fig:ADSB_Louvain_Leiden_Walktrap_size_number}), and betweenness centrality of major hubs had stabilized at near-normal levels (Figure~\ref{fig:ADSB_max_betweeneess_Pagerank_Closeness_centalities}). By contrast, during the April through June 2020 period when edge formation was at its lowest and the network most fragmented, no new variant of concern established itself in Canada.

Several important caveats apply to this comparison. Omicron's rapid global spread was substantially driven by its own biological properties, including high transmissibility and immune escape from prior immunity, which are independent of network structure \cite{omicron_biology_2021}. The four-day interval between global classification and Canadian detection likely reflects both the reconstructed network and the variant's intrinsic fitness advantages over Delta. Additionally, the travel history of the first Canadian Omicron cases involved Nigeria rather than southern Africa, where the variant was first detected. This illustrates a recurring challenge for border surveillance: importation pathways are not always the most geographically obvious ones. A surveillance system calibrated only to direct connections from known source regions would have missed Nigeria entirely. Effective border health surveillance therefore requires monitoring global network connectivity, not just connections from anticipated source countries. The network structure identifies which origin airports have active routes into Canada regardless of whether those origins are currently flagged by health authorities as high-risk, and this is precisely the kind of intelligence that volume-based surveillance cannot provide.

These caveats notwithstanding, the contrast between the 2020 period of structural compression, during which no new variants of concern were established in Canada despite global viral circulation, and the late 2021 period of structural recovery, during which Omicron arrived within days of global detection, is temporally consistent with the hypothesis that network structure shapes the speed and probability of variant importation.

Taken together, these three cases suggest that the temporal network metrics developed here have potential epidemiological validity beyond structural characterization alone. The periods during which our analysis identifies constrained network connectivity coincide with the observed absence or delay of variant importation into Canada. The periods of network recovery coincide with successful and rapid variant establishment. Future work integrating genomic surveillance data with temporal network analysis at the route level could test these associations more rigorously, particularly by examining whether specific dissolved and reformed edges correspond to the travel histories of documented importation cases.


\subsection{Volume Collapse and Network Compression}

The pandemic's impact on human mobility is starkly visible in the passenger volume data. From January through March 2020, the air transportation network operated at near-normal levels. Beginning in April 2020, network conditions changed substantially. Passenger volumes collapsed to just 22\% of their April 2019 levels, a 78\% decline in a single month (Figure~\ref{fig:ADSB_world_volumes}). This is consistent with Canadian border statistics \cite{StatisticsCanada_2026}: Canada experienced a 92\% drop in traveller arrivals (from 7.7 million in April 2019 to 614,000 in April 2020). The synchronous collapse across both measurement systems confirms this was not a localized phenomenon: the entire North American-European air network compressed simultaneously.

Fewer passengers means fewer infected travellers crossing borders. If a pathogen's importation risk is proportional to travel volume, then April 2020 represented a 78--92\% reduction in the structural opportunity for disease introduction. This compression persisted far longer than initially expected. By December 2021, 20 months after the initial collapse, passenger volumes were still only 46\% of pre-pandemic levels (based on Canadian border data: 2.97 million arrivals versus the pre-pandemic baseline of 6--8 million monthly). A pathogen emerging in January 2020 would have faced vastly different transmission landscapes in April 2020 versus April 2021, even though both were pandemic years.

The network did not simply reduce in volume; it lost routes sharply. The number of active airports (nodes) fell from a summer 2019 peak of approximately 9,600 to 8,900 in April 2020, a 7\% reduction (Figure~\ref{fig:ADSB_Number_of_Vertices}). More striking was the collapse in routes: the number of edges fell from a peak of 135,000 in August 2019 to just 80,000 in April 2020, a 40\% drop from peak (Figure~\ref{fig:ADSB_Number_of_Edges}). The winter 2019 baseline was already around 86,000--99,000 edges, so April 2020 represented an approximate return to winter minimum levels, but occurring in spring when traffic would normally be rising, making the disruption particularly visible against seasonal expectations.

\begin{figure}[htbp]
 \center
    \includegraphics[width=\textwidth]{Figures/Air-Transportation/summed_vol_month_plot.pdf}
    \caption{Total number of travellers in the ADS-B air transportation network during 2019--2021.}
    \label{fig:ADSB_world_volumes}
\end{figure}

\begin{figure}[htbp]
    \begin{subfigure}[b]{0.5\textwidth}
        \includegraphics[width=\textwidth]{Figures/Air-Transportation/Directed_Network_number_of_edges.pdf}
        \caption{Evolution of the number of edges.}
        \label{fig:ADSB_Number_of_Edges}
    \end{subfigure}
    \begin{subfigure}[b]{0.5\textwidth}
        \includegraphics[width=\textwidth]{Figures/Air-Transportation/Directed_Network_number_of_nodes.pdf}
        \caption{Evolution of the number of nodes.}
        \label{fig:ADSB_Number_of_Vertices}
    \end{subfigure}
    \caption{Temporal analysis of network structures in the ADS-B air transportation network. The left panel shows the number of edges (routes); the right panel shows the number of nodes (airports).}
    \label{fig:ADSB_network_evolution}
\end{figure}

A 41\% reduction in routes means entire transmission pathways simply ceased to exist. A flight from Toronto to Frankfurt that operated weekly in March 2020 might not have operated at all in April 2020, directly reducing the probability that an infected traveller in one region could reach another. The network contracted dramatically, cutting off connections that would normally accelerate disease spread.

The recovery of nodes and edges was rapid and ultimately exceeded pre-pandemic levels, and this asymmetry with passenger volume is the paper's central finding: \textit{structural recovery not only preceded volumetric recovery but substantially surpassed the pre-pandemic baseline while passenger volumes remained depressed.} Airports came back online quickly; by May--June 2020, node counts had already returned to and exceeded the pre-pandemic peak of approximately 9,600, reaching over 10,000 by mid-2020. Routes followed: by October 2020 edge counts exceeded 150,000, above the pre-pandemic maximum of 135,000. Throughout 2021, both nodes and edges remained consistently above pre-pandemic levels, with edges reaching approximately 180,000--200,000, roughly 60--70\% above the pre-pandemic average of 112,000. This means the network was not merely recovering: it was actively expanding beyond its pre-pandemic scale. This distinction is precisely why passenger volume monitoring alone is insufficient: the structural pathways through which disease moves had not only reopened but multiplied, while the volume of travellers using them remained at roughly half of pre-pandemic levels. Importation risk tracked structural scale, not volumetric recovery.


\subsection{Fragmentation and Community Structure}

The networks are not strongly connected (Figure~\ref{fig:ADSB_number_strong_components_directed_undirected}). Before the pandemic, the number of strongly connected components ranged from approximately 2,100 to 2,750 per month, reflecting seasonal variation in network connectivity. April 2020 produced 2,396 strongly connected components, within the pre-pandemic range, not above it. The more notable structural shift came in May 2020, when SCC dropped to approximately 2,000, falling below the pre-pandemic minimum. This represents the clearest network-level signature of community severance: the global air transportation system, which had functioned as a single highly connected super-network, fragmented into a larger number of smaller, locally-contained clusters. The long-haul inter-regional routes that stitched continents together dissolved first, leaving regional sub-networks that were internally connected but increasingly cut off from one another. From a public health perspective, this is the mechanism through which travel restrictions achieved containment: not simply reducing the volume of movement, but breaking the structural pathways through which a pathogen could traverse from one region to another in a small number of hops. By mid-2020, SCC had stabilised near 2,000--2,100 and remained in that range through 2021, somewhat below the pre-pandemic average of approximately 2,300.

\begin{figure}[htbp]
    \begin{subfigure}[b]{\textwidth}
    \includegraphics[width=\textwidth]{Figures/Air-Transportation/Directed_Network_strongly_connected_components.pdf}
    \caption{Directed networks.}
    \label{fig:ADSB_number_components_strong_directed}
    \end{subfigure}
    \begin{subfigure}[b]{\textwidth}
    \includegraphics[width=\textwidth]{Figures/Air-Transportation/Undirected_Network_strongly_connected_components.pdf}
    \caption{Undirected networks.}
    \label{fig:ADSB_number_components_undirected}
    \end{subfigure}
    \caption{Monthly count of strongly connected components in directed and undirected ADS-B air transportation networks.}\label{fig:ADSB_number_strong_components_directed_undirected}
\end{figure}

From a public health standpoint, fragmentation is a double-edged phenomenon. Greater fragmentation means disease cannot flow as easily between isolated sub-networks; a pathogen circulating in one region might not reach another for weeks or months because the routes no longer exist.

Our edge formation data identifies when this vulnerability window opened: by mid-2021, route formation was running substantially above pre-pandemic levels, and by late 2021, community structure had reverted toward pre-pandemic configurations across both detection algorithms. Consistent with the temporal alignment documented in Section~\ref{subsec:variant_alignment}, the timing of three successive variant importation events into Canada (Alpha in December 2020, Delta in summer 2021, and Omicron in November 2021) corresponds to these successive stages of network recovery. We do not claim that network reconstruction caused these importation events; variant biology, traveller behaviour, and public health policy all contributed independently. However, the correspondence between the period of maximum network compression (April through September 2020), during which no new variant of concern established itself in Canada, and the periods of network recovery, during which each successive variant arrived with decreasing delay, suggests that structural network conditions formed part of the landscape within which importation risk evolved.

Communities represent semi-isolated subnetworks within which disease can spread efficiently, but between which transmission may be more limited (Figure~\ref{fig:ADSB_Louvain_Leiden_Walktrap_size_number}). We used two algorithms (Louvain and Infomap) to detect communities and assess whether the air network remained globally connected or fragmented into regional clusters. Each algorithm identifies clusters differently: Louvain optimizes for modularity, while Infomap uses information-theoretic principles.

\begin{figure}[htbp]
  \begin{subfigure}[b]{0.5\textwidth}
    \includegraphics[width=\textwidth]{Figures/Air-Transportation/communities_louvain_number_timeseries.pdf}
    \caption{Number of communities (Louvain).}
    \label{fig:ADSB_Louvain_number_com}
  \end{subfigure}
  \hfill
  \begin{subfigure}[b]{0.5\textwidth}
    \includegraphics[width=\textwidth]{Figures/Air-Transportation/communities_louvain_max_timeseries.pdf}
    \caption{Size of the maximum community (Louvain).}
    \label{fig:ADSB_Louvain_sie_max_com}
  \end{subfigure}
  \hfill
    \begin{subfigure}[b]{0.5\textwidth}
    \includegraphics[width=\textwidth]{Figures/Air-Transportation/communities_infomap_number_timeseries.pdf}
    \caption{Number of communities (Infomap).}
    \label{fig:ADSB_Walktrap_number_com}
  \end{subfigure}
  \hfill
  \begin{subfigure}[b]{0.5\textwidth}
    \includegraphics[width=\textwidth]{Figures/Air-Transportation/communities_infomap_max_timeseries.pdf}
    \caption{Size of the maximum community (Infomap).}
    \label{fig:ADSB_Walktrap_sie_max_com}
  \end{subfigure}
  \caption{Evolution of the number of communities and size of the maximum community in the ADS-B air transportation network, using Louvain and Infomap community detection algorithms.}
  \label{fig:ADSB_Louvain_Leiden_Walktrap_size_number}
\end{figure}

Before the pandemic, Louvain detected between 19 and 47 communities per month, reflecting high natural variability in how the algorithm partitioned the network across seasons. April 2020 produced 47 Louvain communities, at the upper end of the pre-pandemic range but not beyond it. The number of communities alone therefore does not unambiguously signal pandemic-driven fragmentation. By contrast, the modularity score does: Louvain $Q$ rose from a pre-pandemic range of approximately 0.62--0.65 to a peak of approximately 0.69 in April--May 2020, indicating that while the number of communities did not dramatically increase, the communities that existed became more internally dense and more isolated from one another. Infomap community counts (which operate on a different scale, ranging 200--330 pre-pandemic) followed a similar directional pattern, with $Q$ peaking at approximately 0.62 in the same window. By late 2021, Louvain community counts had converged to 14--16 per month, lower and more stable than the volatile 2019 baseline, and modularity values had declined toward pre-pandemic levels, indicating that the network had reintegrated into a more uniformly connected structure.

A key question for both algorithms is whether the pandemic actually increased network modularity, meaning the degree to which communities became more internally dense and inter-community connections more sparse. Modularity score ($Q$) quantifies this: values near 0 indicate weak community structure, while values approaching 1 indicate strong partitioning. In pre-pandemic months, Louvain modularity scores ranged from approximately $Q \approx 0.62$--$0.65$, consistent with a geographically structured but globally integrated network. By April--May 2020, modularity peaked at approximately $Q \approx 0.69$, indicating that pandemic-era route cancellations genuinely increased the isolation between communities even without dramatically changing their number. Infomap $Q$ followed the same pattern, rising from a pre-pandemic baseline of approximately 0.57--0.60 to a peak near 0.62 in the same window. This finding is directly relevant to the concept of \textit{community severance} in the transport and health literature: travel restrictions in 2020 did not merely reduce movement volume, they structurally tightened regional clusters and weakened the inter-community connections that enable rapid global disease dispersal. The gradual return to lower modularity values through 2021 tracks the reopening of these cross-community pathways and corresponds temporally to the period of accelerated variant importation documented in Section~\ref{subsec:variant_alignment}. Both algorithms agreed on the directional pattern: peak modularity in spring 2020 and gradual decline through 2021, though they differ in the absolute community counts they produce, reflecting their different optimization objectives.

\begin{figure}[htbp]
    \centering
    \includegraphics[width=\textwidth]{Figures/Air-Transportation/modularity_timeseries.pdf}
    \caption{Monthly modularity scores ($Q$) for the ADS-B air transportation network, computed using Louvain and Infomap community detection algorithms. Both algorithms show a peak in April--May 2020, coinciding with the period of maximum route dissolution and network fragmentation, followed by a gradual decline through 2021 as cross-community connectivity was restored.}
    \label{fig:ADSB_modularity_timeseries}
\end{figure}

If disease spreads more easily within communities but is limited between them, then a highly fragmented network containing disease within regional boundaries is beneficial from a public health perspective. April 2020 provided exactly this: the network fractured into more isolated communities, reducing the probability of rapid global spread.

By late 2021, community structure began reverting to pre-pandemic patterns. As communities merged back together, new transmission pathways opened. Recovery involved not only restoring passenger volumes but reinstating the global connectivity that enables rapid pathogen dispersal. Understanding community structure can therefore inform targeted interventions: travel restrictions between communities may be more effective than restrictions within them, and the temporal changes documented here identify specific periods when such restrictions would have had greatest impact.

\subsection{Centrality Shifts: Which Airports Control Disease Introduction?}

Betweenness centrality measures how much an airport lies on the shortest paths between other airports (Figure~\ref{fig:ADSB_max_betweenss_centality}). Chicago O'Hare International Airport (KORD) consistently exhibited the highest betweenness centrality among United States airports from 2019 to 2021, reflecting its role as a major national and international hub. Exceptions occurred in specific months: in July 2019 and July 2021, Wittman Regional Airport (KOSH) in Wisconsin displayed the highest centrality, while in May 2021, Teterboro Airport (KTEB) in New Jersey led.

\begin{figure}[htbp]
    \begin{subfigure}[b]{\textwidth}
    \includegraphics[width=\textwidth]{Figures/Air-Transportation/Max_Betweenness Centrality_in_Each_Graph.pdf}
    \caption{Maximum monthly betweenness centrality.}
    \label{fig:ADSB_max_betweenss_centality}
    \end{subfigure}
    \begin{subfigure}[b]{\textwidth}
    \includegraphics[width=\textwidth]{Figures/Air-Transportation/Max_PageRank_in_Each_Graph.pdf}
    \caption{Maximum monthly PageRank centrality.}
    \label{fig:ADSB_max_Pagerank_centality}
    \end{subfigure}
    \caption{Monthly maximum centrality measures for air transportation networks. Subfigure (a) shows maximum betweenness centrality; subfigure (b) shows maximum PageRank centrality. Airport codes above each data point identify the airport with the highest centrality value in each monthly snapshot.}\label{fig:ADSB_max_betweeneess_Pagerank_Closeness_centalities}
\end{figure}

These observations illustrate an important point: airports with relatively small passenger volumes can still exhibit high betweenness centrality if they serve as critical connectors within specific regional transportation networks. Wittman Regional Airport's July spikes reflect its role as host of EAA AirVenture Oshkosh, the world's largest annual aviation gathering.

It is worth addressing directly whether the KOSH spikes represent genuine public health signal or structural artifacts. EAA AirVenture brings hundreds of thousands of visitors to Oshkosh, but the majority arrive on small general aviation aircraft (Cessnas, Pipers, and similar light planes), rather than commercial airliners. Our network edges are weighted by estimated passenger count, assigning a volume of 2 passengers to flights where the aircraft type could not be identified (the majority of general aviation). A network of several thousand single-engine aircraft is structurally different from a network of equivalent node-count served by widebody jets: the former represents high topological connectivity at very low passenger throughput per edge. In the unweighted network, KOSH's structural centrality in July is real; it is genuinely on many shortest paths during the AirVenture period. However, its \textit{epidemiological} importation risk is not proportional to this betweenness score, because the volume of potentially exposed travellers on any single route is an order of magnitude lower than at a commercial hub. This discrepancy is a known limitation of graph-theoretic centrality when applied to heterogeneous networks where edges vary widely in capacity. We therefore recommend that any operational surveillance tier based on centrality metrics incorporate a \textit{passenger-weighted betweenness} measure, where path weights reflect the total passengers that could traverse each edge, rather than purely topological betweenness. Under such a weighting, KOSH's July spikes would diminish substantially relative to KORD. 

From an epidemiological perspective, locations with high betweenness centrality act as critical bridges for disease transmission. An infected individual travelling through a high-betweenness node has a higher probability of seeding infections in multiple downstream locations. Chicago O'Hare's consistent high betweenness centrality suggests its sustained role as a potential disease dissemination hub. Changes in betweenness centrality over time translate directly to changes in importation risk profiles.

PageRank centrality captures something distinct: the quality of an airport's connections (Figure~\ref{fig:ADSB_max_Pagerank_centality}). Addison Airport (KADS) in Texas consistently demonstrated the highest PageRank centrality, with exceptions in July 2019 and July 2021 when Wittman Regional Airport (KOSH) led, August 2021 when Chicago O'Hare led, and November and December 2021 when DeKalb-Peachtree Airport (KPDK) in Georgia showed the highest centrality. KPDK's elevated PageRank in late 2021 is consistent with a pattern observed for other general aviation airports: as the commercial network recovered and hub-to-hub routes were restored, well-connected regional airports serving dense metropolitan areas temporarily exhibited outsized structural importance relative to their passenger volume.

In disease transmission networks, high PageRank indicates locations that serve as epidemiological amplifiers: nodes that are not only well-connected but connected to other well-connected locations, creating pathways for rapid outbreak propagation. A moderately busy airport connected to major international hubs can seed outbreaks across multiple secondary hubs simultaneously, effectively shortening the time to global dispersal. This cascading effect is particularly consequential for diseases with significant asymptomatic transmission, where infected travellers may transit through hub airports undetected. These network bottlenecks become strategic priorities for surveillance and intervention, as disrupting transmission at high-PageRank nodes can disproportionately reduce outbreak spread relative to lower-ranked locations with similar traffic volume.


\subsection{Edge Dynamics: Pathways Opening and Closing}

Edge formation and dissolution reveal the rate at which the network restructures. In April 2020, we observed a dramatic drop in edge formation: far fewer new routes were opened compared to April 2019 and April 2021 (Figure~\ref{fig:ADSB_edges_formed}). Simultaneously, we observed substantial edge dissolution (Figure~\ref{fig:ADSB_edges_dissolved}). Within weeks, the network had contracted to essential connections.

\begin{figure}[htbp]
    \centering
    \begin{subfigure}[b]{\textwidth}
        \centering
        \includegraphics[width=\textwidth]{Figures/Air-Transportation/edges_dissolved.pdf}
        \caption{Edges Dissolved in the Network.}
        \label{fig:ADSB_edges_dissolved}
    \end{subfigure}
    \hfill
    \begin{subfigure}[b]{\textwidth}
        \centering
        \includegraphics[width=\textwidth]{Figures/Air-Transportation/edges_formed.pdf}
        \caption{Edges Formed in the Network.}
        \label{fig:ADSB_edges_formed}
    \end{subfigure}
    \caption{Monthly dynamics of the ADS-B air transportation network. (a) Edges dissolved each month from 2019 to 2021. (b) Edges formed each month over the same period.}
    \label{fig:ADSB_edges_comparison}
\end{figure}

Edge dissolution is the mechanism through which travel restrictions work. Route cancellation directly eliminates transmission pathways; an infected traveller cannot traverse a route that no longer operates. The rapid edge dissolution in spring 2020 directly reduced importation risk by removing pathways faster than the virus could traverse them. The notable decrease in edge dissolutions in May 2020 (compared to May 2019 and May 2021) indicates the network had already contracted to its minimal viable size.


A counterpoint is that edge formation remained sluggish through all of 2020 and much of 2021. When travel began recovering, it initially restored the most critical connections (major hub-to-hub routes) while leaving peripheral connections dormant. This created a period where disease was more likely to flow through major hubs and less likely to reach smaller communities via direct routes, not through the absence of travel, but rather the concentration of movement through a reduced set of pathways.


\subsection{Recovery and Persistent Vulnerability}\label{subsec_recovery_Persistent}

By December 2021, the network had structurally exceeded pre-pandemic capacity: node counts reached approximately 11,000 (versus a pre-pandemic peak of 9,600) and edge counts approached 190,000 (versus a pre-pandemic peak of 135,000 and average of 112,000). Community counts had converged to 14--16 per month, lower and more stable than the volatile 2019 baseline, and modularity had declined from its 2020 peak, indicating a more uniformly integrated network. However, Canadian border arrival data presents a strikingly different picture: arrivals remained 58\% below pre-pandemic baseline (2,966,051 arrivals versus 7,142,000 in December 2019). This divergence between structural expansion and volumetric depression is the clearest evidence that passenger volume monitoring and network structure monitoring capture different things. The network had grown beyond its pre-pandemic scale while the travellers using it remained scarce. A surveillance system calibrated only to passenger volume would have seen continued depression of risk into late 2021. A system monitoring network structure would have identified an expanding, increasingly connected network that was creating more importation pathways even as each individual pathway carried fewer travellers.

The data tells a coherent epidemiological story spanning 2019--2021. From January through March 2020, the network was vulnerable: a newly emerged pathogen could spread globally via intact routes. In April 2020, that window closed rapidly. The network collapsed across every measured dimension: volume (78\% drop), routes (41\% drop), connected airports (12\% drop), and connectivity (increased fragmentation). This compression persisted through much of 2020 and into 2021.

By late 2021, the network was reassembling. Edges reformed, communities reconnected, and centrality measures returned toward pre-pandemic values. The structure recovered faster than the volume of travellers. This matters because a reconstructed network operating at 50\% capacity is epidemiologically different from a reconstructed network at full capacity: the same number of infected travellers reaches fewer total people.

The persistent gap between structural recovery and volumetric recovery in late 2021 suggests that disease importation risk evolved along multiple timescales: highest in pre-pandemic 2019 (full network, no immunity), lowest in April--June 2020 (structural constraint), elevated but shifting through 2020--2021 (recovering structure, low volume), and normalizing by late 2021 (recovering structure and volume). These temporal changes provide epidemiologists with quantifiable measures of when disease importation risk was highest, demonstrating that this risk was not static across the pandemic but shifted substantially as the network recovered.


\section{Discussion}

Our analysis reveals a limitation of surveillance approaches based on passenger volume alone: network structure matters as much as the number of travellers. When Canada experienced a 92\% collapse in arrivals in April 2020, this represented an obvious reduction in importation risk. Of equal importance, however, was the simultaneous restructuring of the network itself: routes disappeared selectively, communities fragmented, and hub airports remained bottlenecks even as the overall system contracted. These structural changes are not captured by monitoring total passenger numbers.

This distinction has real implications for disease importation. If a pathogen's risk is purely proportional to travel volume, then monitoring passenger counts is sufficient. However, our findings suggest this relationship is more complex. In April 2020, the network lost 41\% of its routes while losing only 12\% of its airports. This asymmetry created a situation where infection imported through a hub airport such as Toronto Pearson would have had difficulty reaching smaller communities directly, but easier access through the remaining routes. A surveillance system monitoring only passenger volume would have registered a 92\% reduction in risk. A system accounting for network structure would have identified a different risk profile: lower overall, but concentrated through fewer pathways.

A further definitional precision is warranted. Throughout this paper, ``importation risk'' refers to the structural potential for an infected traveler to enter and seed infection in a new region. This is not equivalent to a node's betweenness centrality alone. A high-betweenness airport may have many transiting passengers who remain in a sterile international transit zone and never enter the host country's population; in which case the epidemiological seeding potential is substantially lower than the topological centrality score implies. Our analysis does not distinguish between transit and entry passengers, as this information is not available in the ADS-B dataset. This is a consequential limitation: KORD's high betweenness includes a large share of connecting passengers who do not clear customs in Chicago. 

This matters for future pandemics. In the event of novel pathogen emergence, passenger counts convey one dimension of importation risk; network structure analysis conveys another, and both are needed for comprehensive assessment.

Static network analyses are fundamentally limited during crises. A snapshot of the air transportation network taken in April 2020 would show a different structure than one taken in May 2020, and different again in December 2020. Each snapshot accurately represents that moment, but none captures the ongoing process: a system actively restructuring. Our 36-month dataset makes this evolution visible. The number of communities increased in spring 2020, indicating fragmentation into regional clusters. By late 2021, communities were merging again as routes reopened. Centrality measures shifted as different airports became critical connectors at different times. A static analysis would miss this continuous restructuring entirely.

Disease importation risk is not static during a crisis; it evolves as the network evolves. April 2020 presented one risk profile: low volume and highly fragmented structure. December 2021 presented another: low volume but increasingly integrated structure. The risk in late 2021 differed from April 2020 not only because fewer people were travelling, but because they were travelling through a fundamentally different network topology. Temporal analysis reveals these changes; static snapshots obscure them.

This insight extends beyond COVID-19. Any public health crisis affecting travel (whether travel bans, fuel shortages, infrastructure damage, or political instability) will reshape transportation networks. Understanding how that reshaping affects disease risk requires tracking the network over time, not measuring it once.

Chicago O'Hare remained the most structurally important hub throughout the pandemic as measured by betweenness centrality, consistent with its established role as a major national and international connector. However, PageRank centrality reveals greater complexity: this measure showed more volatility, with different airports becoming critical at different times. Addison Airport in Texas demonstrated high PageRank in several months due to the quality of its connections to other hubs. Wittman Regional Airport spiked in July due to EAA AirVenture.

The practical implication is that focusing surveillance on the largest airports by passenger volume is reasonable but insufficient. Chicago O'Hare matters; so does Toronto Pearson for Canadian importation. However, smaller hubs that are structurally well-connected can also be critical entry points. The pandemic illustrated this: even as passenger volume collapsed, remaining travellers often concentrated through major hubs. If infection entered at a hub, the downstream reach was substantial.

For policy, this suggests a \textit{Centrality-Weighted Surveillance Tier} system. Major international hubs (Toronto Pearson, Vancouver International, Chicago O'Hare) should receive permanent high-intensity screening proportional to their consistently dominant betweenness and PageRank values. A second tier of regionally important secondary airports should receive elevated monitoring during periods when their centrality scores rise above defined thresholds. A third tier of smaller regional airports, like Addison (KADS) or Wittman (KOSH), would trigger temporary high-intensity screening during specific high-centrality windows: KOSH in July each year during EAA AirVenture, or any regional airport that transiently becomes a critical connector as hub-to-hub routes shift. Under this framework, monthly network analysis would update each airport's tier assignment dynamically, so that a new international route opening at a regional airport automatically triggers a reassessment of its screening intensity. This approach directs surveillance resources in proportion to actual structural risk rather than historical passenger volume, and would have flagged the late-2021 recovery of international connectivity as a period warranting heightened monitoring, precisely when Omicron arrived.

Effective pandemic response requires understanding transportation networks in real time. Border agencies tracked arrivals, which was valuable. However, they operated largely independently of network analysis. Had agencies known that in late 2021 international arrivals were recovering slowly while network structure was rapidly reconnecting, they could have anticipated that reopening routes were creating unexpected importation pathways.

More broadly, public health agencies require network intelligence alongside traditional epidemiological surveillance. Knowing that a pathogen is circulating in Europe is one piece of information. Knowing which airports connect Europe to North America, and how those connections are changing, is another that is not typically part of border health screening.

The methods developed here are not specific to Canada or to COVID-19. The OpenSky Network provides global ADS-B coverage, and the community detection and centrality approaches are standard graph-theoretic tools. Any public health authority with access to flight data could replicate this analysis for their own region. The Canada--US--Europe focus serves as a proof of concept with well-validated outcome data, but the framework applies equally to monitoring importation risk in other regions. Freely available data and open-source tools substantially lower the barrier to adoption, requiring no proprietary data sources.

These findings carry concrete implications for public health practice. Most fundamentally, passenger volume monitoring alone is insufficient for importation risk assessment. Public health agencies require complementary network intelligence: access to flight data, capacity to construct and update network representations on a monthly or weekly basis, and automated flagging of structural changes such as rapid edge dissolution, community fragmentation, or centrality shifts. The tools developed here, and the freely available OpenSky data on which they rely, make this feasible at modest cost.

Surveillance resources should be allocated in proportion to network centrality rather than passenger volume alone. Betweenness and PageRank centrality identify airports whose structural position makes them disproportionately important for disease transmission, sometimes independently of their size. For Canada, this means intensive screening at major international hubs such as Toronto Pearson and Vancouver International, but also dynamic monitoring of secondary airports whose centrality shifts as the network evolves. A new international route opening at a regional airport changes that airport's epidemiological importance immediately; network analysis captures this in real time whereas passenger volume statistics lag substantially.

The asymmetric recovery documented here has a direct operational implication: network structure and passenger volume diverged sharply in 2021. By December 2021, structural metrics (nodes, edges, community connectivity) had not only recovered but substantially exceeded pre-pandemic levels. International passenger arrivals, however, remained 58\% below the equivalent 2019 month. A surveillance system calibrated to structural recovery would have correctly identified an expanding network with growing importation pathways. A system calibrated to passenger volume would have continued to register depressed risk. Both signals are real; neither alone is sufficient.

Finally, these findings argue for pandemic preparedness planning that explicitly accounts for network restructuring, not only volume reduction. Networks do not simply shrink during crises; they reorganize, concentrating flow through fewer pathways and amplifying the epidemiological importance of major hubs. Contingency plans should anticipate this concentration effect and ensure that hub airports receive surveillance resources commensurate with their increased structural importance during network compression, not their pre-crisis passenger volumes. Building this capacity requires integrating network science expertise into public health institutions, a gap that COVID-19 has made urgent.

\subsection{Limitations}

This analysis has several constraints that should be acknowledged. The geographic scope is limited to Canada, the United States, and Europe, and our findings cannot be directly extrapolated to importation dynamics in other regions. Passenger volume estimates based on aircraft capacity represent upper bounds; actual numbers of infected travellers could be lower if flights were lightly booked.

A specific and important limitation concerns ADS-B sensor coverage density. The OpenSky Network relies on a volunteer receiver network whose density was substantially higher in North America and Europe than in other regions during 2019--2021. In regions with sparse coverage, individual flights may not have been detected by any receiver, and routes that remained operational could appear dissolved in our network when they were in fact still active. This geographic sensor bias has two consequences. First, our network may overstate the degree of route dissolution in regions with historically lower ADS-B coverage; dissolution events in those regions could reflect loss of sensor coverage rather than actual route cancellation. Second, our validation against Statistics Canada international arrival data is appropriate for the geographic scope of the study, but extrapolating the analysis to regions with weaker ADS-B infrastructure would require independent validation against local aviation authority data. For Canada, the United States, and most of Europe, where ADS-B carriage is legally mandated and sensor coverage is dense, this concern is substantially mitigated; the risk is more acute for routes connecting to regions outside our study area. Future work applying this framework globally should incorporate OpenSky sensor density maps to weight edge observations by coverage confidence and distinguish between sensor-driven and policy-driven route dissolution.

The ADS-B system has coverage gaps in regions where equipment is not mandated, which may affect completeness for routes connecting to those areas. Community detection relied on two algorithms that sometimes produced different results, reflecting the inherent ambiguity of network community structure. Finally, we measured network structure during one specific crisis with one specific pathogen. Whether these findings generalize to other pathogens, other crises, or other transportation networks remains to be established. COVID-19's rapid spread and severity drove dramatic policy responses; a slower-moving pathogen or a different type of crisis may produce different patterns of network restructuring.

\renewcommand\bibname{References}
\phantomsection \addcontentsline{toc}{chapter}{References}
\bibliographystyle{plain20}
\bibliography{ref,epidemio,biblio_Arino_Julien,biblio,math_ecology}

\end{document}